\section{Introduction}

Transformers \cite{vaswani2017attention} underpin the success of modern large language models (LLMs), demonstrating impressive performance across a diverse set of tasks \cite{brown2020language}. A key emergent capability of these models is \textit{in-context learning} (ICL), which allows them to perform novel tasks specified by examples in the input prompt, without updating model parameters \cite{brown2020language, dong2022survey, liu2023pre}. This stands in sharp contrast to standard \textit{in-weights learning} (IWL), where knowledge is gradually encoded into the model parameters during training.
Research into the nature of ICL is extensive, ranging from mechanistic accounts including induction heads \cite{olsson2022context} and function vectors \cite{todd2023function}, to theoretical frameworks that interpret ICL as a form of implicit optimization or Bayesian inference \cite{von2023transformers, dai2022can, xie2021explanation}. Most relevant to this work, however, is the finding that the emergence of ICL relies on specific statistical properties of the training data \cite{chan2022data}.

While the emergence of ICL can be productively understood in terms of rational adaptation to environmental pressures \cite{wurgaft2025context}, this perspective has limitations. By focusing on what the asymptotically optimal strategy is (i.e., ICL or IWL), it struggles to account for the learning dynamics, particularly the factors that precipitate transitions between strategies. This gap is significant because empirical studies show that ICL can be transient, emerging early in training only to diminish or be superseded by IWL as the model converges \cite{singh2023transient}.

In this paper, we argue that evolutionary theory offers a powerful explanatory lens for understanding these learning dynamics. Specifically, we examine the environmental factors that drive transitions between two analogous biological strategies: \textit{plasticity} and \textit{evolution}. Plasticity is the capacity of a single genotype to produce different phenotypes (e.g., morphological or behavioral changes) in response to environmental cues. It allows for adaptation within an individual's lifetime, mirroring how ICL modulates inference based on a prompt. Conversely, evolution involves the slow modification of the genotype across generations via selection to produce specific adaptive traits, analogous to updating model parameters via IWL. Furthermore, just as evolution gives rise to plasticity, the capacity for ICL is acquired through the process of IWL.

Evolutionary theory suggests that the reliability of information across timescales determines which adaptive strategy is favored \cite{stearns1989evolutionary}. Plasticity is generally selected for in fluctuating environments, provided reliable cues are available to guide the adaptive response \cite{stearns1989evolutionary, stephens1989variance}. However, plasticity tends to be selected against when the environment is stable enough to obviate the need for flexibility, when reliable cues are unavailable, or when the inherent costs of plasticity are prohibitive \cite{stearns1989evolutionary, stephens1989variance, dewitt1998costs}. Furthermore, \textit{genetic assimilation} describes the evolutionary process by which plastic traits become fixed, eventually rendering their expression insensitive to the cues that originally produced them \cite{waddington1953genetic, pigliucci2006phenotypic}. This phenomenon parallels the empirically observed transience of ICL, where the model's reliance on the prompt diminishes as the task is consolidated into weights \cite{singh2023transient}.

We hypothesize that analogous factors govern the competition between ICL and IWL in Transformers, determining which strategy dominates at different stages of training. To test this, we operationalize two key dimensions: \textit{cue reliability}, defined as the sufficiency of information within a single prompt to specify the target task, and \textit{environmental stability}, defined as the invariance of the target task across different prompts encountered during training. We manipulate these dimensions in two controlled settings: parametric sinusoid regression and few-shot Omniglot classification. Our results confirm that the evolutionary principle regarding environmental predictability across timescales accurately predicts the learning dynamics of Transformers. This suggests that an ecological view of training environments can offer principled guidance for developing new training methodologies.

% Transformers \cite{vaswani2017attention} underpin the success of modern large language models (LLMs), demonstrating impressive performance across diverse tasks \cite{brown2020language}. A key emergent capability of these models is in-context learning (ICL), where models adapt to novel tasks specified by examples given in the input prompt, without requiring gradient-based parameter updates \cite{brown2020language, dong2022survey, liu2023pre}. This rapid and flexible adaptation contrasts with standard in-weights learning (IWL), where knowledge is gradually encoded into model parameters during training \cite{chan2022data, singh2023transient}. The mechanisms enabling ICL are an active area of research, with theories encompassing induction heads that copy and complete patterns \cite{olsson2022context, elhage2021mathematical}, function vector heads that compute task representations \cite{todd2023function, hendel2023context}, and interpretations of ICL as implicit optimization or Bayesian inference within the forward pass \cite{von2023Transformers, dai2022can, xie2021explanation}.

% While IWL traditionally yields stable knowledge from consistent task exposure, the interplay and developmental trajectory between IWL and ICL remain incompletely understood. Notably, ICL capabilities can be transient, emerging early in training but sometimes diminishing or being superseded by IWL as tasks become more predictable or frequently encountered \cite{singh2023transient, anand2024dual}. Characteristics of the data distribution, such as input noise and dynamic input-label mappings --- which influence task predictability --- have also been shown to affect the balance between ICL and IWL \cite{chan2022data, reddy2023mechanistic}. Despite these advances, a comprehensive framework for why and when Transformers shift their reliance between these learning modes is still developing.

% In this paper, we argue that evolutionary biology offers one explanatory lens for understanding this dynamic interplay. Organisms adapt to variable environments via analogous mechanisms: phenotypic plasticity and genetic evolution. Phenotypic plasticity --- the ability of a genotype to produce different phenotypes (e.g., changes in morphology, physiology, or behavior like learning) in response to environmental cues \cite{west2003developmental, scheiner1993genetics} --- allows rapid, within-lifetime adaptation, akin to a Transformer's ICL adapting to a prompt. This contrasts with genetic evolution --- the slower, generational modification of the genotype --- which encodes more permanent traits, analogous to IWL.

% Evolutionary biology has established that environmental predictability critically determines the favored adaptive strategy \cite{stearns1989evolutionary}. Phenotypic plasticity is generally favored when the environment fluctuates and reliable cues are available to guide an appropriate adaptive response \cite{stearns1989evolutionary, stephens1989variance}. However, plasticity may be selected against if the environment is extremely stable (negating the need for flexibility and favoring genetically hardcoded traits), if cues are unreliable, or if the inherent costs of plasticity are too high \cite{stearns1989evolutionary, stephens1989variance, dewitt1998costs}. Furthermore, a key long-term dynamic is genetic assimilation: if an environment that previously favored a plastic response becomes stable, the consistently advantageous phenotype, initially induced plastically by cues, can become genetically encoded, shifting the trait from flexible to fixed \cite{waddington1953genetic, pigliucci2006phenotypic}. This offers a compelling parallel to the observed transience of ICL.

% Building on this evolutionary framework, we hypothesize that analogous dimensions of task predictability govern which modes of learning Transformers adopt and how their balance shifts over time. We operationalize two key dimensions, inspired by these biological concepts: \textit{cue reliability}, defined as the clarity and informativeness of information within a single prompt for specifying the required task, and \textit{environmental stability}, defined as the consistency or predictability of the task structure across different prompts encountered over the course of training. We predict that high environmental stability will favor IWL, while high cue reliability will enhance ICL efficacy, particularly when stability is low. Moreover, we anticipate that conditions of high stability for tasks initially learned via ICL may lead to an ``assimilation'' of these solutions into IWL.

% To test these hypotheses, we employ a controlled experimental design that systematically investigates the independent and interactive effects of cue reliability and environmental stability on the use and development of ICL and IWL. We deploy this approach in two distinct domains: a parametric sinusoid regression task and a few-shot binary classification task on Omniglot images. Our findings aim to establish predictability as a critical factor shaping adaptive learning in Transformers, offering insights for understanding ICL and potentially guiding future training methodologies.

% Transformer architectures \cite{vaswani2017attention} underpin the success of modern large language models (LLMs), demonstrating impressive performance across diverse tasks \cite{brown2020language}. A key emergent capability in these models is in-context learning (ICL), where models adapt to novel tasks specified by examples within the input prompt, without requiring gradient-based parameter updates \cite{brown2020language, dong2022survey, liu2023pre}. This rapid, flexible adaptation contrasts with standard in-weights learning (IWL), where knowledge is gradually encoded into model parameters during training \cite{chan2022data, singh2023transient}. While previous work has begun to elucidate the specific conditions that govern models' preference between ICL and IWL \cite{chan2022data, singh2023transient, bietti_birth_2023, anand2024dual, he2024context, chan_toward_2025, singh_strategy_2025}, this dynamic interplay nonetheless remains incompletely understood.

% In this paper, we argue that one way to understand this interplay is to look to evolutionary biology, where organisms adapt to variable environments via analogous mechanisms. Phenotypic plasticity---the ability of a genotype to produce different phenotypes in response to environmental cues via processes like learning \cite{west2003developmental, pigliucci2001phenotypic}---allows rapid adaptation within a lifetime, akin to ICL adapting to a prompt. For example, certain animals change the thickness and color of their coats in preparation for changing seasons. This phenotypic change is triggered by number of daylight hours, a cue that is predictive of the relevant environmental conditions (temperature and snowfall) \cite{underwood_photoperiod_1980, otte_snow_2024}. This rapid adaptability contrasts with genetic evolution—the slower, generational modification of the genotype—which encodes more permanent traits (analogous to IWL).

% When do organisms evolve to display phenotypic plasticity? Previous research in evolutionary biology has established that it depends critically on the predictability of the environment \cite{scheiner2021loss, leung2020reduced}. Phenotypic plasticity is generally favored when the environment undergoes fluctuations and reliable cues are available to guide an appropriate adaptive response \cite{scheiner2021loss, leung2020reduced}. However, plasticity may be selected against if the environment is either extremely stable (negating the need for flexibility) or if its variations are too erratic and cues too unreliable to permit effective plastic adjustments \cite{leung_reduced_2020, reed2010evolutionary, tufto2015genetic, chevin2015estimating, bonamour2019phenotypic, stamps2020information, lehtonen2021sensitive, king_evolution_2019}.

% Based on this evolutionary framework, we propose that analogous dimensions of task predictability should govern which modes of learning Transformers adopt. We operationalize cue reliability as the clarity and reliability of information within a single prompt for specifying the required task \cite{dall2015learning, lehtonen2021sensitive, stamps2020information} and 
% environmental stability as the stability or predictability of the task across different prompts over time \cite{boyce2006predicting, leung2020evolution}.
% To test the hypothesis that these factors determine what mode of learning Transformers will adopt, we use a controlled experimental design that systematically investigates the independent and interactive effects of cue reliability and environmental stability on the use of ICL and IWL. We deploy this approach in two distinct domains: a parametric sinusoid regression task and a few-shot binary classification task on Omniglot images.
